% ============================================================
%  Função Quadrática — Apostila Premium | PratiqueJá
%  Engine: XeLaTeX
% ============================================================
\documentclass[12pt]{article}
\usepackage{../../pratiqueja}
\usepackage{tikz}

% \hlm — destaque de resultado em modo-math (cor sem bold)
\newcommand{\hlm}[1]{{\color{darkblue}#1}}

\setsubject{Função Quadrática}

\begin{document}
\pagenumbering{arabic}
\setstretch{1.2}
\color{bodytext}

\heroheader{Função Quadrática}%
           {Parábola, Vértice, Zeros, Sinal e Otimização}%
           {Matemática · Avançado}

\setlength{\columnsep}{18pt}
\setlength{\columnseprule}{0.5pt}
\def\columnseprulecolor{\color{exborder}}

\begin{multicols}{2}

%─────────────────────────────────────────────────────────────
\TW{Def}{badgepurple}{Função Quadrática}{$f(x)=ax^2+bx+c$, com $a\neq0$}

\regra{$f(x) = ax^2 + bx + c$ com $a,b,c\in\mathbb{R}$ e $\mathbf{a\neq0}$.
Domínio $\mathbb{R}$; o gráfico é uma \textbf{parábola}.}

\nota{A \emph{equação} $ax^2+bx+c=0$ pergunta \emph{quando} $f(x)=0$;
a \emph{função} descreve $f(x)$ para \textbf{todo} $x$.}

%─────────────────────────────────────────────────────────────
\TW{Graf}{darkblue}{Gráfico — A Parábola}{Concavidade e eixo de simetria}

\regra{O sinal de $a$ dá a \textbf{concavidade}: $a>0$ abre para
\textbf{cima} (mínimo no vértice); $a<0$ para \textbf{baixo} (máximo).}

\begin{center}
\begin{tikzpicture}[x=0.55cm,y=0.45cm,font=\scriptsize,>=stealth]
\draw[graycolor,dashed,line width=0.5pt] (1,-2.4)--(1,3.0);
\draw[->,line width=0.7pt] (-2.4,0)--(4.6,0) node[right]{$x$};
\draw[->,line width=0.7pt] (0,-2.6)--(0,3.3) node[above]{$y$};
\draw[darkblue,line width=1.2pt,domain=-1.95:3.95,samples=60,smooth] plot (\x,{0.5*((\x-1)^2)-2});
\fill[vered] (-1,0) circle (2.2pt);
\node[vered,anchor=south east] at (-1,0.1){$x_1$};
\fill[vered] (3,0) circle (2.2pt);
\node[vered,anchor=south west] at (3,0.1){$x_2$};
\fill[badgegreen] (1,-2) circle (2.2pt);
\node[badgegreen,anchor=north] at (1.15,-2.05){$V$};
\node[graycolor,anchor=south,font=\scriptsize] at (1,3.0){eixo $x_v$};
\end{tikzpicture}
\end{center}

\exemp{%
  \textbf{$\Delta>0$:} cruza o eixo $x$ em $x_1, x_2$\\[2pt]
  \textbf{$\Delta=0$:} tangencia — raiz dupla $x_v$\\[2pt]
  \textbf{$\Delta<0$:} não cruza o eixo%
}
\obs{\textbf{Eixo de simetria:} reta $x = x_v = \dfrac{-b}{2a}$; as raízes
são simétricas em relação a ele.}

%─────────────────────────────────────────────────────────────
\TW{V}{badgegreen}{Vértice da Parábola}{Ponto de máximo ou mínimo}

\formula{$x_v = \dfrac{-b}{2a} \qquad y_v = \dfrac{-\Delta}{4a}$}

\exemp{%
  \textbf{$f(x)=x^2-10x+9$:}\\
  $\Delta = 100-36 = 64$\\
  $x_v = \dfrac{10}{2} = 5 \qquad y_v = \dfrac{-64}{4} = -16$\\
  \hlm{$V(5,-16)$} — mínimo ($a>0$)%
}

%─────────────────────────────────────────────────────────────
\TW{f(x){=}0}{badgepurple}{Zeros da Função}{Onde a parábola cruza o eixo $x$}

\regra{Os \textbf{zeros} são as raízes de $ax^2+bx+c=0$ (Bhaskara) — as
interseções com o eixo $x$.}
\exemp{%
  $\Delta>0$: dois zeros distintos\\[2pt]
  $\Delta=0$: zero duplo $x_v$ \quad $\Delta<0$: nenhum zero real%
}
\regra{\textbf{Forma fatorada} ($\Delta\geq0$): $\ f(x)=a(x-x_1)(x-x_2)$.}
\exemp{%
  \textbf{$f(x)=2x^2-8x+6$:} $\Delta=16$, $x_1=1$, $x_2=3$\\
  \hlm{$f(x)=2(x-1)(x-3)$}%
}

%─────────────────────────────────────────────────────────────
\TW{$\pm$}{darkblue}{Estudo do Sinal}{Quando $f$ é positiva, nula ou negativa}

\regra{Com $\Delta>0$ e zeros $x_1<x_2$:}
\exemp{%
  \textbf{$a>0$:} $f>0$ fora de $[x_1,x_2]$; \ $f<0$ \emph{entre} as raízes\\[2pt]
  \textbf{$a<0$:} $f<0$ fora; \ $f>0$ \emph{entre} as raízes%
}
\obs{Se $\Delta<0$: $f$ tem \textbf{sinal constante} (igual ao de $a$)
para todo $x$.}

%─────────────────────────────────────────────────────────────
\TW{D/Im}{badgegreen}{Domínio, Imagem e Variação}{Imagem e crescimento pelo vértice}

\regra{$D(f)=\mathbb{R}$. A imagem é limitada por $y_v$:}
\formula{$\text{Im}(f) = \begin{cases} [y_v,+\infty) & \text{se } a>0\\ (-\infty,y_v] & \text{se } a<0 \end{cases}$}
\obs{\textbf{$a>0$:} decresce em $(-\infty,x_v]$, cresce em $[x_v,+\infty)$.
\textbf{$a<0$:} o inverso.}
\exemp{%
  \textbf{$f(x)=x^2-4x+3$:} $x_v=2$, $y_v=-1$\\
  $D=\mathbb{R}$, $\text{Im}=[-1,+\infty)$\\
  \hlm{decresce em $(-\infty,2]$, cresce em $[2,+\infty)$}%
}

\columnbreak

%─────────────────────────────────────────────────────────────
\TW{Opt}{badgepurple}{Aplicações de Otimização}{Máximo e mínimo em situações reais}

\regra{O \textbf{vértice} dá o valor ótimo. Identifique $a,b,c$ e calcule
$x_v$ e $y_v$.}

\SH{P1 — altura de projétil}
\exemp{%
  $h(t)=-5t^2+20t$ ($a<0$ → máximo):\\
  $t_v = \dfrac{-20}{-10} = 2$\,s; \ $h_{\max} = -20+40 = \hlm{20\text{ m}}$%
}
\SH{P2 — maximizar área}
\exemp{%
  Cerca $40$\,m: largura $x$, comprimento $20-x$:\\
  $A(x)=x(20-x)=-x^2+20x$; \ $x_v=10$\\
  $A_{\max} = \hlm{100\text{ m}^2}$%
}
\SH{P3 — lucro máximo}
\exemp{%
  $L(x)=-2x^2+80x-600$: \ $x_v=\dfrac{-80}{-4}=20$\\
  $L_{\max} = -800+1600-600 = \hlm{\text{R\$}\,200}$%
}

\end{multicols}

\end{document}
